\label{ch:data-release}

\section{Publications}

\begin{refbox}
\begin{itemize}
  \item \textbf{Primary source.} B. Aharmim et al. (SNO Collaboration),
    \emph{Electron Energy Spectra, Fluxes, and Day-Night Asymmetries of
    $^8$B Solar Neutrinos from the 391-Day Salt Phase SNO Data Set}, Phys.
    Rev. C \textbf{72}, 055502 (2005), arXiv:nucl-ex/0502021
    \parencite{Aharmim2005SaltPhase}. Appendix A (its Tables XXX-XXXIV) is
    the sole source of every real table this project's
    \texttt{data/detector/sno\_ii/} tree consumes.
  \item \textbf{Earlier flux-only salt-phase result.} S. N. Ahmed et al.
    (SNO Collaboration), \emph{Measurement of the Total Active $^8$B Solar
    Neutrino Flux at the Sudbury Neutrino Observatory with Enhanced Neutral
    Current Sensitivity}, Phys. Rev. Lett. \textbf{92}, 181301 (2004)
    \parencite{Ahmed2004SaltFlux}. Not this document's data source (a
    different, earlier 254-day, flux-only extraction with no day/night
    spectral information), but its own published methodology note is used
    in \cref{sec:data-provenance} as an independent cross-check of this
    document's correlated-covariance construction.
  \item \textbf{Phase-I discovery paper.} Q. R. Ahmad et al. (SNO
    Collaboration), \emph{Direct Evidence for Neutrino Flavor
    Transformation from Neutral-Current Interactions in the Sudbury
    Neutrino Observatory}, Phys. Rev. Lett. \textbf{89}, 011301 (2002)
    \parencite{Ahmad2002NCDiscovery}. Source of the NC-measures-total-active-flux
    principle used throughout \cref{ch:detector-implementation}.
  \item \textbf{External comparison: NuFIT.} I. Esteban et al., \emph{The
    fate of hints: updated global analysis of three-flavor neutrino
    oscillations}, JHEP \textbf{12}, 216 (2024), arXiv:2410.05380
    \parencite{EstebanNuFit2024} -- base analysis publication for NuFIT 6.1
    (2025), \url{www.nu-fit.org}. Source of the starting $(\theta_{12},
    \Delta m^2_{21},\theta_{13})$ values used throughout
    \cref{ch:detector-implementation,ch:results}.
\end{itemize}
\end{refbox}

\section{Data provenance and the transcription process}
\label{sec:data-provenance}

\begin{theorybox}
Unlike Daya Bay's own machine-readable Zenodo/GitHub data release, no
digitised data file exists for \cite{Aharmim2005SaltPhase}: the collaboration
published a dedicated numerical-recipe note
(\texttt{sno.phy.queensu.ca/sno/results\_09\_03/howto.pdf}) alongside their
\emph{earlier}, flux-only salt result \parencite{Ahmed2004SaltFlux}, but no
equivalent machine-readable release for this document's own 391-day
spectral+day/night source. Every real number this document's inference
uses -- the 17-bin CC spectrum, the integrated NC/ES fluxes, the two
$19\times19$ statistical correlation matrices, the systematic percent
sensitivities, and the zenith-livetime exposure -- was therefore transcribed
by hand from the rendered PDF of \cite{Aharmim2005SaltPhase}'s own Appendix A
into \texttt{data/detector/sno\_ii/}, with every transcription independently
cross-checked wherever the paper's own internal redundancy allowed it.
\end{theorybox}

\begin{resultbox}
Four independent cross-checks passed, all recorded in
\texttt{data/detector/sno\_ii/metadata/source.json}:
\begin{enumerate}
  \item The 17 transcribed CC day bins sum to $1.7325\times10^6\,{\rm
    cm^{-2}s^{-1}}$, matching the paper's own independently quoted day
    integral flux, $1.73\times10^6\,{\rm cm^{-2}s^{-1}}$ (night: $1.6422$
    vs.\ $1.64$).
  \item Both $19\times19$ statistical correlation matrices have an exact
    unit diagonal, are symmetric (after resolving one genuine
    $10^{-4}$-level rounding inconsistency in the printed table itself,
    verified against the rendered PDF twice), and are positive-semi-definite.
  \item The 60-bin zenith-exposure table's bin edges match the paper's own
    worked example (``the top bin has $\cos\theta_z=-0.5$ to $-0.4667$'')
    exactly, and its total livetime sums to 391.43 days -- matching the
    paper's own ``391-Day Salt Phase'' title/dataset name to within
    rounding.
  \item The salt-phase response-function coefficients quoted in this
    project's own design notes were verified to match the primary source's
    Eq.~A3 exactly, resolving an initial ambiguity with a different,
    similarly-worded formula from the earlier flux-only paper's own
    methodology note \parencite{Ahmed2004SaltFlux}.
\end{enumerate}
\end{resultbox}

\begin{notebox}
\textbf{Two entries could not be transcribed with confidence, and are
disclosed rather than guessed.} Table XXXIV's \texttt{vertex\_resolution}
and \texttt{angular\_resolution} rows are missing a value partway through
their own CC$_9$-CC$_{17}$ block in the rendered source; two independent
transcription attempts produced misaligned columns rather than a consistent
reading. Both rows' CC$_9$-CC$_{17}$ entries are set to \texttt{0.0} (an
explicit placeholder, not a measurement) in
\texttt{cc\_spectral\_systematics.csv}, flagged in
\texttt{metadata/source.json}, and excluded from this document's own
systematic-covariance discussion of dominant contributions
(\cref{ch:results}). Separately, Tables XIX/XX (NC's own two dedicated
systematics: internal neutron background, neutron capture efficiency) were
never transcribed at all -- the systematic covariance used throughout this
document therefore understates NC's true systematic uncertainty specifically.
\end{notebox}

\section{Table catalogue}

\Cref{tab:snoii-tables-observation,tab:snoii-tables-covariance} list every
real table this project's \texttt{tpeanuts.detector.sno\_ii.io} module
loads, and the physics each one encodes, developed in full in
\cref{ch:detector-implementation}.

\newcommand{\snoiicell}[1]{\parbox[t]{0.42\textwidth}{\raggedright #1}}

\begin{table}[htb]
\centering
\caption{Real published observation tables (Tables XXX, XXIV of the primary source).}
\label{tab:snoii-tables-observation}
\begin{tabular}{@{}l l l@{}}
\toprule
\textbf{Cached file} & \textbf{Real contents} & \textbf{Loader} \\
\midrule
\code{cc\_spectrum\_day\_night.csv} & \snoiicell{The 17-bin CC equivalent-flux spectrum (Table XXX), day and night, statistical uncertainty.} & \pyfunc{load\_cc\_spectrum\_day\_night} \\[4pt]
\code{integrated\_fluxes\_day\_night.csv} & \snoiicell{Day/night NC and ES equivalent fluxes, stat+syst (Table XXIV); the CC integral is also carried, as a cross-check only.} & \pyfunc{load\_integrated\_fluxes\_day\_night} \\[4pt]
\code{zenith\_livetime.csv} & \snoiicell{60-bin real $\cos(\text{zenith})$ livetime distribution, 391.43-day total (Table XXXI).} & \pyfunc{load\_zenith\_exposure} \\
\bottomrule
\end{tabular}
\end{table}

\begin{table}[htb]
\centering
\caption{Real published covariance and systematics tables (Tables XXXII-XXXIV).}
\label{tab:snoii-tables-covariance}
\begin{tabular}{@{}l l l@{}}
\toprule
\textbf{Cached file} & \textbf{Real contents} & \textbf{Loader} \\
\midrule
\code{statistical\_correlation\_\{day,night\}.csv} & \snoiicell{The two real $19\times19$ statistical correlation matrices from SNO's own signal extraction.} & \pyfunc{load\_statistical\_correlation} \\[4pt]
\code{cc\_spectral\_systematics.csv} & \snoiicell{16 detector systematics $\times$ 19 channels, signed percent sensitivity (Table XXXIV).} & \pyfunc{load\_cc\_spectral\_systematics} \\
\bottomrule
\end{tabular}
\end{table}

\begin{notebox}
Not fetched/reproduced: Tables XIX/XX's NC-specific systematics (see above);
Table XVII's fundamental constants and radiative-correction factors
($\omega_{\rm CC}$, $\omega_{\rm NC}$, $\omega_{\rm ES}$); and the earlier
Phase-I and 254-day flux-only salt datasets, which are already covered by
this project's separate, sibling \texttt{detector.sno} package (Phase I) and
are not combined with this document's own salt-phase-only fit.
\end{notebox}
